\documentclass[a4paper,11pt]{jsarticle}
\usepackage{amsmath}

\title{確認問題}
\author{学籍番号: 24RB1074 \and 氏名: 半田悠人}
\date{\today}

\begin{document}
\maketitle

\section*{問題1：3手法の使い分け}

\begin{enumerate}
    \item \textbf{50都市のTSP、CPU時間1分} \\
    \textbf{解答：} 焼きなまし法 (SA) \\
    \textbf{理由：} 50都市のTSPは組合せ最適化問題であり、局所探索だけでは局所最適解から抜け出せない可能性がある。SAでは、探索の初期に悪化する解も確率的に受け入れるため、限られた計算時間内でも局所最適解を抜け出し、比較的良い解を探索できる。
    \item \textbf{1000個の品物のナップサック問題} \\
    \textbf{解答：} 遺伝的アルゴリズム (GA) \\
    \textbf{理由：} 1000個の品物から選択する大規模な組合せ問題であり、解候補を0と1の列として表現しやすい。GAでは多数の解候補を用いて広い範囲を探索できるため、与えられた3手法の中では適している。

    \item \textbf{解空間が小さく、滑らかな問題} \\
    \textbf{解答：} 局所探索 \\
    \textbf{理由：} 局所最適解に陥る危険が小さいため、実装が簡単で計算コストの低い局所探索でも、良好な解を効率よく得られると期待できる。

    \item \textbf{並列計算機を使える環境での大規模スケジューリング問題} \\
    \textbf{解答：} 遺伝的アルゴリズム (GA) \\
    \textbf{理由：} GAでは複数の個体の評価を独立に行えるため、並列計算を利用しやすい。また、多数の解候補を同時に探索するので、大規模なスケジューリング問題にも適用しやすい。
\end{enumerate}

\section*{問題2：SAの挙動}

焼きなまし法の受容確率の公式 $P = \exp(-\Delta E/T)$ に基づく。ここでは、現在の解より目的関数値が悪化する量を $\Delta E>0$ とする。

\begin{enumerate}
    \item \textbf{$\Delta E=5$、$T=10$ のとき、受容確率 $P$ を計算せよ} \\
    \textbf{解答：} $P \approx 0.607$ \\
    \textbf{解説：} 公式に代入すると、以下のようになる。
    \begin{equation*}
        P = \exp\left(-\frac{5}{10}\right) = \exp(-0.5) \approx 0.607
    \end{equation*}
    したがって、受容確率は約60.7\%である。

    \item \textbf{同じ $\Delta E$ で $T=1$ のとき、$P$ はどう変わるか} \\
    \textbf{解答：} $P \approx 0.0067$ となり、受容確率は大幅に小さくなる。 \\
    \textbf{解説：} 公式に代入すると、以下のようになる。
    \begin{equation*}
        P = \exp\left(-\frac{5}{1}\right) = \exp(-5) \approx 0.0067
    \end{equation*}
    したがって、受容確率は約0.67\%まで低下し、悪化する解をほとんど受け入れなくなる。

    \item \textbf{なぜ温度を徐々に下げるのか自分の言葉で説明せよ} \\
    \textbf{解答例：} 最初から温度が低いと、少しでも悪くなる解をほとんど選ばなくなり、途中で見つけた解から抜け出せなくなる可能性がある。そのため、最初は温度を高くしてさまざまな解を試し、そこから少しずつ温度を下げて、より良い解に絞っていくのである。
\end{enumerate}

\end{document}
